\section{Training Methods}


\subsection{Supervised Fine-Tuning (SFT)}

\subsubsection{Full Fine-Tuning}

Full fine-tuning updates all model parameters via standard next-token prediction:

\begin{equation}
\mathcal{L}_{\text{SFT}} = -\sum_{i=1}^{T} \log P_\theta(y_i | y_{<i}, x)
\end{equation}

where $x$ is the input, $y$ is the target sequence, and $\theta$ represents all trainable parameters.

\textbf{Advantages}: Maximum flexibility, can adapt to very different distributions.

\textbf{Disadvantages}: Memory-intensive (requires optimizer states for all parameters), risk of catastrophic forgetting.

\subsubsection{Low-Rank Adaptation (LoRA)}

LoRA~\citep{hu2021lora} freezes pre-trained weights and injects trainable low-rank matrices:

\begin{equation}
h = W_0 x + \frac{\alpha}{r} BA x
\end{equation}

where $W_0 \in \mathbb{R}^{d \times k}$ is frozen, $B \in \mathbb{R}^{d \times r}$ and $A \in \mathbb{R}^{r \times k}$ with $r \ll \min(d, k)$.

Gym implements several LoRA variants:

\begin{itemize}
\item \textbf{Standard LoRA}: Rank $r = 8-64$, $\alpha = 16-32$.
\item \textbf{rsLoRA}~\citep{kalajdzievski2023rslora}: Rank-stabilized scaling $\alpha/\sqrt{r}$.
\item \textbf{DoRA}~\citep{liu2024dora}: Weight-decomposed adaptation improving magnitude/direction optimization.
\item \textbf{PiSSA}~\citep{meng2024pissa}: Principal singular value/vector initialization for faster convergence.
\item \textbf{LongLoRA}~\citep{chen2023longlora}: Shifted sparse attention for context extension.
\end{itemize}

\subsubsection{Quantized LoRA (QLoRA)}

QLoRA~\citep{dettmers2023qlora} combines 4-bit quantization with LoRA:

\begin{equation}
W_{\text{quant}} = \text{Quantize}(W_0, \text{NF4}) + BA
\end{equation}

Using NormalFloat4 (NF4) quantization and double quantization for quantization constants, QLoRA achieves:
\begin{itemize}
\item 4x memory reduction vs FP16 base model
\item Minimal quality degradation (< 1\% on MMLU)
\item Enables 65B training on single 48GB GPU
\end{itemize}

Gym's QLoRA implementation includes:
\begin{itemize}
\item Automatic bit-width selection (4/8-bit)
\item Double quantization for constants
\item Paged AdamW optimizer for memory spikes
\item Gradient checkpointing integration
\end{itemize}

\subsection{Preference-Based Methods}

\subsubsection{Direct Preference Optimization (DPO)}

DPO~\citep{rafailov2023dpo} eliminates the reward model by directly optimizing policy from preference data:

\begin{equation}
\mathcal{L}_{\text{DPO}} = -\mathbb{E}_{(x, y_w, y_l)} \left[ \log \sigma \left( \beta \log \frac{\pi_\theta(y_w|x)}{\pi_{\text{ref}}(y_w|x)} - \beta \log \frac{\pi_\theta(y_l|x)}{\pi_{\text{ref}}(y_l|x)} \right) \right]
\end{equation}

where $y_w$ is the preferred output, $y_l$ is the dispreferred output, $\pi_{\text{ref}}$ is the reference policy, and $\beta$ controls deviation.

Gym supports:
\begin{itemize}
\item Standard DPO with frozen reference model
\item Identity Preference Optimization (IPO) variant
\item Conservative DPO (cDPO) for distribution shift
\end{itemize}

\subsubsection{Kahneman-Tversky Optimization (KTO)}

KTO~\citep{ethayarajh2024kto} optimizes from binary feedback without paired comparisons:

\begin{equation}
\mathcal{L}_{\text{KTO}} = \mathbb{E}_{(x,y,z)} \left[ z \cdot v(x,y) - (1-z) \cdot v(x,y) \right]
\end{equation}

where $z \in \{0, 1\}$ indicates thumbs-up/down and $v(x,y)$ is the value function.

\textbf{Advantage}: Requires only binary labels, not ranked pairs, reducing annotation cost by 50-75\%.

\subsubsection{Odds Ratio Preference Optimization (ORPO)}

ORPO~\citep{hong2024orpo} combines SFT and preference learning in a single stage:

\begin{equation}
\mathcal{L}_{\text{ORPO}} = \mathcal{L}_{\text{SFT}} + \lambda \mathbb{E} \left[ \log \sigma \left( \log \frac{P(y_w|x)}{P(y_l|x)} \right) \right]
\end{equation}

\textbf{Advantage}: Removes need for separate SFT stage, reducing total training time by 30-40\%.

\subsection{Reinforcement Learning Methods}

\subsubsection{Proximal Policy Optimization (PPO)}

PPO~\citep{schulman2017ppo} is the classical RLHF method requiring separate reward model:

\begin{equation}
\mathcal{L}_{\text{PPO}} = \mathbb{E}_t \left[ \min(r_t(\theta) A_t, \text{clip}(r_t(\theta), 1-\epsilon, 1+\epsilon) A_t) \right]
\end{equation}

where $r_t(\theta) = \frac{\pi_\theta(a_t|s_t)}{\pi_{\text{old}}(a_t|s_t)}$ is the importance ratio, $A_t$ is the advantage, and $\epsilon$ is the clipping parameter.

Gym's PPO implementation includes:
\begin{itemize}
\item Value head for advantage estimation
\item Generalized Advantage Estimation (GAE)
\item KL penalty term for stability
\item Support for external reward models
\end{itemize}

\subsubsection{Group Relative Policy Optimization (GRPO)}

GRPO~\citep{shao2024deepseekmath} eliminates the value network by computing advantages from group statistics:

\begin{equation}
A_i^g = r_i - \frac{1}{G} \sum_{j=1}^G r_j
\end{equation}

where group $g$ contains $G$ rollouts for the same query, and advantages are relative to group mean.

\textbf{Algorithm}:
\begin{enumerate}
\item Generate $G$ responses per query (typically $G=8$)
\item Compute reward $r_i$ for each response
\item Calculate group-relative advantages
\item Apply PPO-style clipped objective
\end{enumerate}

\textbf{Advantages over PPO}:
\begin{itemize}
\item No value network required (reduces memory 30-40\%)
\item More stable advantages (invariant to reward scale)
\item Better sample efficiency (DeepSeek-Math achieved 52.4\% on MATH benchmark)
\end{itemize}

Gym's GRPO trainer supports:
\begin{itemize}
\item Configurable group size $G = 1-16$
\item Optional advantage normalization
\item Ground truth filtering for homogeneous groups
\item Multi-turn dialogue rollouts
\end{itemize}

\subsubsection{Training-Free GRPO}

Our most significant contribution is Training-Free GRPO~\citep{tencent2025trainingfree}, which operates entirely in context space:

\textbf{Core Idea}: Instead of updating model parameters via gradients, maintain a library of semantic experiences (natural language insights) injected into context.

\textbf{Three-Stage Process}:

\textbf{Stage 1 - Trajectory Summarization}: For each rollout $(q, o_i, r_i)$:
\begin{verbatim}
Prompt: "Summarize this trajectory step-by-step:
1. Which experiences were used?
2. Where did errors occur?
3. What was the outcome?"
\end{verbatim}

\textbf{Stage 2 - Group Advantage Extraction}: For group of $G$ summaries:
\begin{verbatim}
Prompt: "Compare successful vs failed trajectories.
Suggest updates: Add/Modify/Delete experiences.
Focus on strategic patterns, max 32 words each."
\end{verbatim}

\textbf{Stage 3 - Batch Consolidation}: Across all groups in batch:
\begin{verbatim}
Prompt: "Consolidate suggested updates.
Merge similar experiences, ensure no duplication.
Output: JSON operations [{"option": "merge", ...}]"
\end{verbatim}

\textbf{Algorithm}:
\begin{verbatim}
E = {}  # Experience library
for epoch in [1..3]:
    for batch in dataset:
        for query in batch:
            # Generate G rollouts with experiences injected
            outputs = [π(o|q, E.format_for_prompt())
                      for _ in range(G)]
            rewards = [R(q, o) for o in outputs]

            # Skip homogeneous groups
            if std(rewards) < threshold:
                continue

            # Extract semantic advantages
            summaries = [LLM.summarize(q, o, r)
                        for o, r in zip(outputs, rewards)]
            operations = LLM.extract_insights(summaries, E)

        # Batch consolidation
        final_ops = LLM.consolidate(all_operations, E)
        E.apply_operations(final_ops)

    # Model parameters θ unchanged!
\end{verbatim}

\textbf{Experience Format Example}:
\begin{verbatim}
[G0]. When solving geometry problems with intersections,
      validate solutions lie within bounded regions, not
      extensions, to avoid extraneous answers.

[G10]. When using mathematical invariants to prove
       impossibility, validate against known achievable
       states or small test cases.
\end{verbatim}

\textbf{Results on AIME Benchmarks}:
\begin{table}[h]
\centering
\begin{tabular}{lcc}
\toprule
Method & AIME24 & AIME25 \\
\midrule
Vanilla GRPO & 80.0\% & 67.9\% \\
Training-Free GRPO & \textbf{82.7\%} & \textbf{73.3\%} \\
\midrule
Training Cost & \$10,000+ & \$18 \\
Training Samples & 10,000+ & 100 \\
Training Time & 20 GPU-days & 6 hours \\
\bottomrule
\end{tabular}
\caption{Training-Free GRPO achieves superior performance at 99.8\% cost reduction.}
\label{tab:grpo-results}
\end{table}

\textbf{Key Insights}:
\begin{enumerate}
\item \textbf{Ground truth helps but optional}: Performance drops 2-4\% without ground truth labels, but self-discrimination via majority voting still improves over baseline.
\item \textbf{Multi-epoch essential}: Single pass insufficient; 3 epochs yields best results.
\item \textbf{Strong base model required}: Works excellently on DeepSeek-V3.1 (671B) and Qwen3-32B, degrades on weaker models.
\item \textbf{Cross-domain transfer}: Frozen model + domain experiences outperforms specialized fine-tuned models:
   \begin{itemize}
   \item ReTool (math-tuned): 67\% AIME, 18\% web tasks
   \item MiroThinker (web-tuned): 43.5\% AIME, 53.6\% web tasks
   \item Training-Free GRPO: 82.7\% AIME, 67.8\% web tasks
   \end{itemize}
\end{enumerate}

\subsection{Generalized Sequential Policy Optimization (GSPO)}

GSPO extends GRPO to sequence-level optimization with policy regularization:

\begin{equation}
\mathcal{L}_{\text{GSPO}} = \mathcal{L}_{\text{GRPO}} + \lambda \text{KL}(\pi_\theta || \pi_{\text{ref}}) + \alpha \mathcal{R}_{\text{stability}}
\end{equation}

Additional terms:
\begin{itemize}
\item KL divergence from reference policy prevents distribution collapse
\item Stability regularizer for Mixture-of-Experts (MoE) models
\end{itemize}

GSPO is particularly effective for:
\begin{itemize}
\item Very large models (30B+ parameters)
\item MoE architectures (Mixtral, DeepSeek-MoE)
\item Long sequences (1K+ tokens)
\end{itemize}

